% TEMPLATE-MILC-v3.tex - plantilla maestra UNICA de las guias MILC v3.
%
% La usan las cuatro familias de guias, cada una con su driver:
%   · Grados 6-11          -> scripts/build-guias-g11.py
%   · Bebras               -> scripts/build-guias-bebras.py
%   · Semillero            -> scripts/build-guias-semillero.py
%   · Territorio interior  -> scripts/build-guias-territorio-interior.py
%
% Antes habia tres copias casi identicas (~400 lineas iguales, 6 distintas):
% cada mejora habia que aplicarla tres veces, y en la practica no se hacia
% (el motor de recursos vivia solo en dos de ellas). Lo que varia entre
% familias esta parametrizado en 6 placeholders que cada driver llena
% (se nombran aqui SIN su sintaxis real: el builder reemplaza por texto plano,
% tambien dentro de los comentarios, y un valor multilinea rompe el preambulo):
%
%   HEADER_IZQ        encabezado izquierdo de pagina
%   HEADER_DER        encabezado derecho de pagina
%   PORTADA_ETIQUETA  rotulo sobre el numero grande (GRADO/MOMENTO/MODULO)
%   PORTADA_SERIE     linea de serie de la portada
%   PORTADA_ID        identificador de la guia en la portada
%   PORTADA_CIRCULO   el circulito de grado (°), o vacio si no aplica
%   PDF_TITULO        titulo en texto plano (sin \textbf ni comillas TeX) para
%                     los metadatos del PDF (/Title); lo produce lib_xelatex.texto_pdf
%
% Composicion (auditoria 2026-09-03): las bandas de estacion piden espacio
% (\Needspace*) y prohiben el salto justo despues (\nopagebreak), asi no quedan
% huerfanas al pie; las cajas largas (softbox, drawbox, infoband, puentebox) son
% `breakable`, asi una caja de media pagina ya no arrastra todo a la siguiente
% dejando la anterior al 40 %. Todo texto sobre mostaza o verde va en milcNegro
% (blanco sobre #E5B400 da 1,93:1 y sobre #58A923 2,95:1; ninguno pasa WCAG AA).
% El resto de claves las documenta content/guias/_SCHEMA.md. El script valida
% que no queden placeholders sin reemplazar antes de escribir el archivo final.

% Etiquetado PDF (latex-lab): /Lang, estructura y texto alternativo de las
% figuras, para lectores de pantalla. Debe ir ANTES de \documentclass.
\DocumentMetadata{lang=es-CO,tagging=on}
\documentclass[11pt,letterpaper]{article}
% headheight=24pt: la cabecera derecha lleva dos lineas (identidad de la guia
% y estacion actual). headsep se acorta en la misma medida para que el bloque
% de texto quede exactamente donde estaba con headheight=12pt.
\usepackage[margin=2.0cm,headheight=24pt,headsep=13pt]{geometry}
\usepackage[table]{xcolor}
\usepackage{fontspec}
\usepackage[spanish]{babel}
\usepackage{tikz}
% Librerías TikZ para los diagramas de las guías (bloque `recursos:`):
%  · babel  → babel en español vuelve activos caracteres como «>», y TikZ los usa
%             en sus opciones (>=stealth, ->). Sin esta librería, cualquier
%             diagrama con flechas revienta con «\language@active@arg> has an
%             extra }». Debe ir ANTES de que se dibuje nada.
%  · positioning        → sintaxis «below left=2pt and 3pt».
%  · decorations.pathreplacing → llaves ({brace}) para señalar rangos.
\usetikzlibrary{babel,positioning,decorations.pathreplacing}
\usepackage{graphicx}
% Circuitos: el trabajo de electrónica se hace hoy con micro:bit y sus sensores
% integrados (MakeCode, sin cableado), así que no se necesitan esquemáticos.
% Si algún día entran montajes con protoboard, basta descomentar esta línea:
% los diagramas .tex/.tikz declarados en `recursos:` ya se incrustan solos.
% \usepackage{circuitikz}
\usepackage[most]{tcolorbox}
\usepackage{tabularx}
\usepackage{array}
\usepackage{fancyhdr}
\usepackage{titlesec}
\usepackage[document]{ragged2e}  % bandera a la derecha en todo el cuerpo
\usepackage{enumitem}
\usepackage{microtype}
\usepackage{etoolbox}
\usepackage{needspace}
% hyperref al final del preambulo: metadatos del PDF (titulo, autor, idioma,
% familia), marcadores por estacion (\pdfbookmark) y enlaces sin recuadro.
\usepackage[hidelinks,
  pdftitle={Presentaciones interactivas — caminos que vuelven al menú},
  pdfauthor={PhD. Álvaro Cárdenas Orozco},
  pdfsubject={Tecnología e Informática · Grado 8 · Guía 22 · MILC v3},
  pdflang={es-CO},
  bookmarksopen=true,bookmarksnumbered=false]{hyperref}

% Helvetica Neue no trae la flecha U+2192, asi que cada «→» escrito literalmente
% en el YAML se compilaba a NADA, en silencio (462 flechas en 61 guias: rutas del
% tipo «paleta → tipografia → lockup» salian rotas). Mapeamos el caracter a su
% version matematica, que si tiene glifo. Asi el autor puede seguir escribiendo
% «→» comodamente en el YAML.
\usepackage{newunicodechar}
\newunicodechar{→}{\ensuremath{\rightarrow}}
% Mismo problema con los simbolos de marca y vinetas: el «✓» de «una palomita ✓»
% salia como cuadrito vacio en el PDF de todas las guias que lo usaban.
\usepackage{pifont}
\newunicodechar{✓}{\ding{51}}
\newunicodechar{✗}{\ding{55}}
\newunicodechar{✔}{\ding{52}}
\newunicodechar{•}{\textbullet}
\newunicodechar{≥}{\ensuremath{\geq}}
\newunicodechar{≤}{\ensuremath{\leq}}

\setmainfont{Helvetica Neue}
\setsansfont{Helvetica Neue}
\setlength{\parindent}{0pt}
\setlength{\parskip}{6pt}
% Sin particion de palabras: en 6.º-7.º una palabra cortada por guion al final
% de linea («Comuni-cacion») frena la lectura mucho mas que una linea corta.
\hyphenpenalty=10000
\exhyphenpenalty=10000
\pretolerance=2000
\tolerance=3000
\setlength{\emergencystretch}{3em}
\linespread{1.10}
\renewcommand{\arraystretch}{1.25}
\setlist{leftmargin=5mm,itemsep=1mm,topsep=1mm}
\newcolumntype{Y}{>{\RaggedRight\arraybackslash}X}

\definecolor{milcMagenta}{HTML}{E6007E}
\definecolor{milcVino}{HTML}{5A0038}
\definecolor{milcTurquesa}{HTML}{0093A5}
\definecolor{milcVerde}{HTML}{58A923}
\definecolor{milcMostaza}{HTML}{E5B400}
\definecolor{milcOcre}{HTML}{B86600}
\definecolor{milcNegro}{HTML}{241816}
\definecolor{milcGris}{HTML}{F7F7F4}
\definecolor{milcLinea}{HTML}{E8E2DD}
\definecolor{milcGradePrimary}{HTML}{FF6600}
\definecolor{milcGradeDark}{HTML}{9A3E00}
\definecolor{milcGradeSoft}{HTML}{FFE4D0}
\definecolor{milcGradeAccent}{HTML}{CFE7FF}
\definecolor{milcGradeText}{HTML}{FFFFFF}
\color{milcNegro}

\pagestyle{fancy}
\fancyhf{}
% Cabecera de dos lineas: identidad de la guia arriba y, a la derecha y en
% gris, la estacion en curso (\markright desde \milcestacion). Antes de la
% primera estacion la segunda linea va vacia; el \strut mantiene la altura
% para que la linea de arriba no baile entre paginas.
\lhead{\small Tecnología e Informática · Grado 8\\[-1pt]{\footnotesize\strut}}
\rhead{\small Guía 22 · MILC v3\\[-1pt]{\footnotesize\strut\color{milcNegro!65}\rightmark}}
\cfoot{\small\thepage}
\renewcommand{\headrulewidth}{0pt}

% Sobre mostaza (#E5B400) y verde (#58A923) el texto va en milcNegro; sobre el
% resto de la paleta, en blanco. Se usa dentro de fonttitle/fontupper, donde un
% \color posterior manda sobre coltitle.
\newcommand{\milctintasobre}[1]{%
  \ifboolexpr{test{\ifstrequal{#1}{milcMostaza}} or test{\ifstrequal{#1}{milcVerde}}}%
    {\color{milcNegro}}{\color{white}}}

\newtcolorbox{titlebox}[1]{
  enhanced,colback=#1,colframe=#1,arc=7mm,boxrule=0pt,
  left=7mm,right=7mm,top=3mm,bottom=3mm,
  before skip=8pt,after skip=8pt,
  fontupper=\sffamily\bfseries\Large\color{white}
}

\newtcolorbox{softbox}[3]{
  enhanced,breakable,pad at break=3mm,
  colback=#2,colframe=#1,borderline west={4pt}{0pt}{#1},
  arc=2mm,boxrule=.4pt,left=4mm,right=4mm,top=3mm,bottom=3mm,
  before skip=6pt,after skip=8pt,title={#3},coltitle=white,
  colbacktitle=#1,fonttitle=\sffamily\bfseries\milctintasobre{#1},fontupper=\RaggedRight,
  attach boxed title to top left={xshift=3mm,yshift=-2mm},
  boxed title style={arc=2mm, boxrule=0pt}
}

\newtcolorbox{drawbox}[2]{
  enhanced,breakable,pad at break=4mm,
  colback=white,colframe=#1,arc=4mm,boxrule=1pt,
  left=5mm,right=5mm,top=4mm,bottom=4mm,before skip=8pt,after skip=8pt,
  title={#2},coltitle=white,colbacktitle=#1,fonttitle=\sffamily\bfseries\milctintasobre{#1},
  fontupper=\RaggedRight,
  attach boxed title to top left={xshift=4mm,yshift=-2mm},
  boxed title style={arc=3mm, boxrule=0pt}
}

\newtcolorbox{infoband}[1]{
  enhanced,breakable,pad at break=2mm,colback=#1!8,colframe=#1!50,arc=2mm,boxrule=.5pt,
  left=4mm,right=4mm,top=2mm,bottom=2mm,before skip=4pt,after skip=4pt,
  fontupper=\small\RaggedRight
}

% Puente narrativo entre secciones (contrato editorial Conéctate).
% Da continuidad: "Acabas de X. Ahora vas a Y porque Z."
% Visualmente: banda izquierda en color del grado, texto en itálica pequeña.
\newtcolorbox{puentebox}{
  enhanced,breakable,pad at break=2mm,colback=milcGradeSoft,colframe=milcGradeDark,
  borderline west={3pt}{0pt}{milcGradeDark},
  arc=0mm,boxrule=0pt,left=5mm,right=4mm,top=2.5mm,bottom=2.5mm,
  before skip=4pt,after skip=8pt,
  fontupper=\itshape\small\color{milcNegro}\RaggedRight
}
% La flecha va en modo matematico a proposito: Helvetica Neue NO tiene el glifo
% U+2192, asi que \textrightarrow se compilaba a NADA («Missing character: There
% is no → in font Helvetica Neue») y el puente salia sin flecha. En modo
% matematico la flecha viene de la fuente matematica y si se dibuja.
\newcommand{\puente}[1]{%
  \begin{puentebox}%
    \textcolor{milcGradeDark}{$\rightarrow$}\hspace{2mm}#1%
  \end{puentebox}%
}

\newcommand{\linead}[1]{\noindent\color{milcLinea}\rule{#1}{0.5pt}\color{milcNegro}}
\newcommand{\respuesta}[1]{\par\vspace{2mm}\foreach \n in {1,...,#1}{\noindent\color{milcLinea}\rule{\linewidth}{0.5pt}\par\vspace{5mm}}\color{milcNegro}}
\newcommand{\checkbox}{\textcolor{milcGradeDark}{\fbox{\phantom{x}}}\hspace{5pt}}

% ─── Listas aireadas para las guías socioemocionales ────────────────────
% Componer las verificaciones como «\checkbox texto\\» deja el texto que salta
% de línea debajo del cuadrito y apelmaza el bloque. Estos entornos alinean la
% continuación y airean los ítems. Los usa Territorio Interior; cualquier
% familia puede hacerlo.
\newlist{checklista}{itemize}{1}
\setlist[checklista]{
  label=\textcolor{milcGradeDark}{\fbox{\phantom{x}}},
  leftmargin=9mm, labelsep=3.2mm, itemsep=2.4mm,
  topsep=2.5mm, parsep=0pt, partopsep=0pt, before=\RaggedRight
}
\newlist{pasoslista}{enumerate}{1}
\setlist[pasoslista]{
  label=\textcolor{milcGradeDark}{\sffamily\bfseries\arabic*.},
  leftmargin=8mm, labelsep=3mm, itemsep=2.8mm,
  topsep=1mm, parsep=0pt, partopsep=0pt, before=\RaggedRight
}

\newcommand{\phasecard}[4]{
\begin{tcolorbox}[enhanced,colback=#3,colframe=#2,arc=3mm,boxrule=.5pt,height=3.05cm,valign=center,halign=center,left=1.5mm,right=1.5mm,top=2mm,bottom=2mm]
{\sffamily\bfseries\fontsize{22}{24}\selectfont\textcolor{#2}{#1}}\par
{\sffamily\bfseries\small #4}
\end{tcolorbox}
}

% ── Piezas del rediseno 2026 (legibilidad 6.º-11.º) ───────────────────
% Banda de estacion: reemplaza el titlebox macizo. Titulo a la izquierda,
% dato practico (tiempo/modalidad) a la derecha, en una sola linea.
% Nunca huerfana: pide 9 lineas libres (o salta de pagina rellenando con
% \vfill) y prohibe el salto justo despues de la banda. Nueve y no seis:
% la caja partible que sigue necesita ~80pt (titulo adosado, relleno y dos
% lineas) para arrancar en la misma pagina; con menos, tcolorbox la manda
% entera a la siguiente y la banda se queda sola. El penalty 10000 va
% en `after`, ANTES del espacio posterior: un `after skip` (aunque sea 0pt)
% es un \vskip tras la caja, y ese glue es un punto de corte legal que un
% \nopagebreak escrito despues de \end{tcolorbox} ya no cierra. Cada banda
% deja marcador PDF y actualiza la cabecera derecha.
\newcounter{milcestacion}
\newcommand{\milcestacion}[2]{%
\Needspace*{9\baselineskip}%
\stepcounter{milcestacion}%
\pdfbookmark[1]{#2}{milcestacion\themilcestacion}%
\markright{#2}%
\begin{tcolorbox}[enhanced,colback=#1,colframe=#1,arc=2mm,boxrule=0pt,
  left=5mm,right=5mm,top=2.6mm,bottom=2.6mm,before skip=11pt,
  after={\par\nopagebreak\vskip7pt}]
{\sffamily\bfseries\fontsize{15}{18}\selectfont\milctintasobre{#1}#2}
\end{tcolorbox}}

% Tarjeta de paso numerada: sustituye las tablas «Pilar / Aplicacion», que
% eran formato de consulta docente, no de estudiante. El numero va en un
% circulo sobre el borde para que la secuencia se lea de un vistazo.
\newcommand{\milcpaso}[3]{%
\Needspace{4\baselineskip}%
\begin{tcolorbox}[enhanced,colback=white,colframe=milcLinea,arc=1.5mm,boxrule=.9pt,
  left=14mm,right=4mm,top=3mm,bottom=3mm,before skip=4pt,after skip=4pt,
  fontupper=\RaggedRight,
  overlay={\node[anchor=north west,circle,fill=#2,text=white,inner sep=0pt,
    minimum size=7.4mm,font=\sffamily\bfseries\small]
    at ([xshift=3.6mm,yshift=-3.2mm]frame.north west) {#1};}]
#3
\end{tcolorbox}}

% Casilla marcable de tamano util con lapiz (la anterior era \fbox{\phantom{x}},
% de alto variable segun el texto que la seguia).
\newcommand{\milcmarca}{\framebox[3.8mm]{\rule{0pt}{3.4mm}}}

% Fila marcable de lista suelta (checks de la estacion 1).
\newcommand{\milccheckrow}[1]{%
\par\vspace{1.8mm}\noindent
\makebox[6.4mm][l]{\raisebox{-.55mm}{\textcolor{milcNegro}{\milcmarca}}}%
\parbox[t]{\dimexpr\linewidth-6.4mm\relax}{\RaggedRight #1}\par}

% Celda marcable dentro de la rubrica (tabla de autoevaluacion). Misma
% sangria francesa, pero pensada para vivir dentro de una columna X.
\newcommand{\milcopcion}[1]{%
\makebox[5.2mm][l]{\raisebox{-.55mm}{\textcolor{milcNegro}{\milcmarca}}}%
\parbox[t]{\dimexpr\linewidth-5.2mm\relax}{\RaggedRight #1}}

% ── Recursos visuales de la guía ──────────────────────────────────────
% Declarados en el bloque `recursos:` del YAML y emitidos por el builder.
% \guiaFigura   → imagen rasterizada o PDF (foto, carta celeste, montaje).
% \guiaDiagrama → fuente TikZ/circuitikz (.tex) incrustada como vector:
%                 es la vía para circuitos y esquemáticos editables.
% [#1] = texto alternativo (alt) · #2 = ruta relativa al .tex · #3 = pie
% El alt llega al PDF etiquetado (/Alt) y lo lee el lector de pantalla.
\newcommand{\guiaFigura}[3][]{%
\begin{center}
\includegraphics[width=0.88\linewidth,height=0.38\textheight,keepaspectratio,alt={#1}]{#2}\\[1.5mm]
{\footnotesize\itshape\textcolor{milcNegro!75}{#3}}
\end{center}
}

\newcommand{\guiaDiagrama}[2]{%
\begin{center}
\input{#1}\\[1.5mm]
{\footnotesize\itshape\textcolor{milcNegro!75}{#2}}
\end{center}
}

\begin{document}
\thispagestyle{empty}

% ── Portada ───────────────────────────────────────────────────────────
% Antes: un rectangulo de color a pagina completa. Se veia bien en pantalla,
% pero en la fotocopiadora del colegio se comia el toner y salia gris sucio.
% Ahora la banda ocupa el tercio superior y el resto va en blanco.
\begin{tikzpicture}[remember picture,overlay]
  \path[fill=milcGradePrimary]
    (current page.north west) rectangle ([yshift=-10.6cm]current page.north east);

  \node[anchor=north west,text=milcGradeText] at ([xshift=2.0cm,yshift=-1.55cm]current page.north west)
    {\sffamily\bfseries\fontsize{12}{14}\selectfont GRADO};
  \node[anchor=north west,text=milcGradeText,opacity=.85] at ([xshift=2.0cm,yshift=-2.35cm]current page.north west)
    {\sffamily\bfseries\fontsize{8.5}{10}\selectfont SERIE GUÍAS · TECNOLOGÍA E INFORMÁTICA};
  \node[anchor=north east,text=milcGradeText,opacity=.85,align=right] at ([xshift=-2.0cm,yshift=-1.55cm]current page.north east)
    {\sffamily\bfseries\fontsize{9}{12}\selectfont I.E. Sor María Juliana\\Cartago, Valle del Cauca};

  \node[anchor=west,text=milcGradeText] (gradonum) at ([xshift=1.85cm,yshift=-7.1cm]current page.north west)
    {\sffamily\bfseries\fontsize{112}{112}\selectfont 8};
  % El circulito de grado (°) solo aplica a las guias de grado («11°»); en
  % Bebras y Semillero el numero grande es un momento/modulo, no un grado.
  % Se posiciona relativo al nodo `gradonum`, no en coordenadas absolutas.
  \draw[milcGradeText,line width=1.6pt]
    ([xshift=2.0mm,yshift=-4.5mm]gradonum.north east) circle (.15cm);

  \node[anchor=west,text=milcGradeText,text width=11.4cm,align=left] at ([xshift=1.5cm,yshift=0.5cm]gradonum.east)
    {
     {\sffamily\bfseries\fontsize{9}{11}\selectfont\MakeUppercase{Período 3 · Multimedia, narrativa y ciberseguridad}}\\[3mm]
     {\sffamily\bfseries\fontsize{21}{25}\selectfont\setlength{\baselineskip}{27pt}Caminos\\[4pt]que vuelven al menú}\\[3.5mm]
     {\sffamily\bfseries\fontsize{15}{18}\selectfont Guía 22-8-TIC}
    };
\end{tikzpicture}

\vspace*{9.4cm}
\vfill

{\sffamily\bfseries\fontsize{8}{10}\selectfont\textcolor{milcGradeDark}{LO QUE TE LLEVAS HOY}}

\begin{tcolorbox}[enhanced,colback=white,colframe=milcNegro,arc=2mm,boxrule=1.6pt,
  left=5mm,right=5mm,top=4mm,bottom=4mm,before skip=3pt,after skip=12pt,
  fontupper=\RaggedRight\fontsize{11}{15.5}\selectfont]
Una presentación interactiva en Slides, Canva, Genially o PowerPoint sobre un tema escolar real, con al menos cinco pantallas conectadas por hipervínculos que dejan elegir el camino y volver al menú, tres animaciones que ayudan a entender y no solo decoran, y una nota de cinco líneas que explique por qué esos caminos y no una secuencia.
\end{tcolorbox}

\begin{tcolorbox}[enhanced,colback=milcGris,colframe=milcGris,arc=2mm,boxrule=0pt,
  left=5mm,right=5mm,top=3.5mm,bottom=3.5mm,after skip=12pt]
{\sffamily\bfseries\fontsize{8}{10}\selectfont\textcolor{milcNegro!55}{ESTA GUÍA ES DE}}\\[3mm]
\begin{tabularx}{\linewidth}{X p{3.2cm} p{3.2cm}}
\linead{\linewidth} & \linead{3.0cm} & \linead{3.0cm}\\[-1mm]
{\fontsize{8.5}{10}\selectfont\textcolor{milcNegro!55}{Nombre completo}} &
{\fontsize{8.5}{10}\selectfont\textcolor{milcNegro!55}{Grupo}} &
{\fontsize{8.5}{10}\selectfont\textcolor{milcNegro!55}{Fecha}}\\
\end{tabularx}
\end{tcolorbox}

\vfill

\noindent\textcolor{milcLinea}{\rule{\linewidth}{1pt}}\\[2mm]
{\sffamily\bfseries\fontsize{9.5}{12}\selectfont Autor: Dr. Álvaro Cárdenas Orozco}\\[1mm]
{\sffamily\fontsize{8.5}{11}\selectfont\textcolor{milcNegro!55}{Metodología MILC · apertura ancestral · pensamiento computacional · triángulo Dussel–estoicismo–Floridi}}

\newpage

\section*{Apertura: Presentaciones interactivas --- caminos que vuelven al menú}

\begin{softbox}{milcGradeDark}{milcGradeSoft}{Saber ancestral}
En Guapi y en los ríos de la costa caucana, los \textbf{decimeros} recitan de memoria poemas de 44 versos: las \textbf{décimas glosadas}. La estructura es estricta. Primero va una copla de cuatro versos que «lleva la esencia». Después vienen cuatro décimas de diez versos. Cada décima tiene que terminar en el verso correspondiente de la copla inicial (Oslender, 2003, citando a Pedrosa y Vanín, 1994). Es decir: el poema abre con un menú de cuatro líneas y desarrolla cada una por su lado. Cada desarrollo vuelve al menú antes de pasar al siguiente. Nadie se pierde porque cada rama regresa. Y las décimas no tienen autoría fija: cada decimero añade o cambia partes, y así se vuelven de todos. La \textbf{cara de exclusión}: Oslender llama a estas décimas «discursos ocultos de resistencia». Las comunidades negras del Pacífico las usan para contar lo que la historia oficial calla y para reclamar su territorio. Y los decimeros son mayores a los que cada vez se escucha menos. Hoy vas a construir una presentación como una décima glosada: un menú que lleva la esencia y ramas que siempre vuelven.\par\smallskip{\footnotesize\textcolor{milcNegro!70}{\textbf{Fuente.} Oslender, U. (2003). «Discursos ocultos de resistencia»: tradición oral y cultura política en comunidades negras de la costa pacífica colombiana. \emph{Revista Colombiana de Antropología, 39}, 203--236.}}
\end{softbox}

\begin{softbox}{milcTurquesa}{milcGris}{Saber contemporáneo}
Una \textbf{presentación interactiva} no se recorre en línea recta. Quien la mira elige por dónde ir, entra a una sección, regresa y salta a otra. Se construye con dos herramientas. Los \textbf{hipervínculos internos}: un texto o botón que lleva a otra pantalla; en Slides, clic derecho, «Insertar enlace», y eliges la diapositiva. Las \textbf{animaciones}: elementos que aparecen, se mueven o cambian cuando haces clic. La regla es una sola: cada enlace y cada animación sirve para entender, no para decorar. Si quitas la animación y se entiende igual, sobra. Tres estructuras cubren casi todo: el \textbf{menú}, una pantalla índice con botones a cada sección; el \textbf{glosario}, una palabra que lleva a su definición y devuelve; y el \textbf{retorno}, el botón «volver al menú» que va en todas las pantallas. Sin retorno, quien entra a una rama queda atrapado. Antes de abrir la herramienta, el mapa se dibuja en papel: nodos, flechas y regresos.
\end{softbox}

\begin{softbox}{milcMostaza}{milcGris}{Pregunta-puente}
\textit{En la décima glosada, cada rama termina volviendo a un verso de la copla inicial. Si tu presentación tiene un botón que entra a una sección y ninguno que vuelva, ¿qué le pasa a quien hizo clic?}
\end{softbox}

\begin{softbox}{milcVerde}{milcGris}{Saber y saber hacer}
Al terminar podrás: (1) \textbf{identificar} en un mapa de papel los nodos, las ramas y los regresos de una narrativa; (2) \textbf{crear} una presentación de cinco pantallas con menú, hipervínculos y retorno; (3) \textbf{evaluar} cada animación con la prueba de quitarla; (4) \textbf{explicar} por qué un camino con ramas necesita regresos.
\end{softbox}



\small
\begin{tabularx}{\textwidth}{>{\bfseries}p{3.0cm}Y}
\rowcolor{milcGradePrimary!12}Autor & Dr. Álvaro Cárdenas Orozco\\
DBA & Producción multimedia ética y comportamiento responsable en entornos digitales (MEN, T\&I)\\
Referentes & Dussel (estética de la liberación) · Ley 1336 · Floridi (infoesfera)\\
Producto final & Una presentación interactiva en Slides, Canva, Genially o PowerPoint sobre un tema escolar real, con al menos cinco pantallas conectadas por hipervínculos que dejan elegir el camino y volver al menú, tres animaciones que ayudan a entender y no solo decoran, y una nota de cinco líneas que explique por qué esos caminos y no una secuencia.\\
\end{tabularx}
\normalsize

\puente{En la sesión 1 hiciste que una pieza se entendiera en tres segundos. Hoy conectas varias pantallas para que quien mira elija su camino. En la sesión 3 vas a contar algo en tres minutos con tu voz.}

\milcestacion{milcGradeDark}{Ruta MILC: así vamos a aprender}
\begin{tabularx}{\textwidth}{>{\centering\arraybackslash}X>{\centering\arraybackslash}X>{\centering\arraybackslash}X>{\centering\arraybackslash}X}
\phasecard{1}{milcVerde}{milcGris}{Escucha\\[-1mm]\small Observo, escucho y pregunto.} &
\phasecard{2}{milcTurquesa}{milcGris}{Entender\\[-1mm]\small Sistematizo y ordeno.} &
\phasecard{3}{milcMagenta}{milcGris}{Hacer\\[-1mm]\small Construyo y aplico.} &
\phasecard{4}{milcVino}{milcGris}{Cierre\\[-1mm]\small Evalúo y reflexiono.}
\end{tabularx}


\puente{Antes de abrir la herramienta, vas a dibujar el mapa en papel: cinco nodos, flechas y regresos. Es la copla inicial de la décima, la que lleva la esencia.}

\milcestacion{milcVerde}{Estación 1 de 3 · Escucha}

\begin{softbox}{milcVerde}{milcGris}{Escena para escuchar}
\textbf{Actividad 1 · IDENTIFICA --- El mapa en papel} (15 min · individual). (1) Elige un tema escolar que tenga partes: el ciclo del agua, las reglas de un deporte, cómo funciona el micro:bit, la historia del colegio. (2) Dibuja cinco círculos: uno es el menú y los otros cuatro son secciones. (3) Traza flechas del menú a cada sección. (4) Traza las flechas de regreso, de cada sección al menú, con otro color. (5) Marca con una estrella la sección que necesita una palabra explicada, ahí irá el glosario.
\end{softbox}

{\sffamily\bfseries\fontsize{8}{10}\selectfont\textcolor{milcNegro!55}{MARCA CUANDO LO HAYAS HECHO}}
\milccheckrow{Tengo cinco nodos, uno de ellos es el menú.}
\milccheckrow{Cada sección tiene flecha de ida y flecha de regreso.}
\milccheckrow{Marqué dónde va el glosario.}

\begin{infoband}{milcVerde}


\textbf{Lo que va al cuaderno (Actividad 1 · IDENTIFICA).} \textbf{Título:} «Actividad 1 --- El mapa en papel». \textbf{Formato:} el mapa de cinco nodos con las flechas de ida en un color y las de regreso en otro, y la estrella del glosario. \textbf{Extensión:} media página. \textbf{Sabes que terminaste cuando} desde cualquier nodo puedes seguir una flecha de vuelta al menú con el dedo.
\end{infoband}

\puente{Ya tienes el mapa. Ahora aprendes los tres tipos de enlace y las tres animaciones útiles, y con tu pareja armas el menú y las primeras ramas.}

\milcestacion{milcTurquesa}{Estación 2 de 3 · Entender}

\begin{softbox}{milcTurquesa}{milcGris}{Lo que vas a entender}
\textbf{Actividad 2 · CREA --- El menú y las ramas} (30 min · parejas). (1) Con tu pareja, en Slides o la herramienta que indique tu docente, creen cinco diapositivas y pongan el título de cada una según el mapa. (2) En la primera, el menú, escriban cuatro botones o textos, uno por sección. (3) Seleccionen cada botón, clic derecho, «Insertar enlace», y elijan la diapositiva de esa sección. (4) En cada sección pongan el contenido en pocas palabras, con la jerarquía de la sesión 1. (5) Prueben en modo presentación: desde el menú, entren a las cuatro secciones.
\end{softbox}

\milcpaso{1}{milcTurquesa}{\textbf{Los tres enlaces.} Menú: la pantalla índice con un botón por sección; deja elegir por dónde empezar. Glosario: una palabra subrayada que lleva a su definición y trae de vuelta. Retorno: el botón «volver al menú» en todas las pantallas menos el menú.}
\milcpaso{2}{milcTurquesa}{\textbf{Las tres animaciones útiles.} Revelar: un elemento aparece después de otro, primero la pregunta y luego la respuesta. Resaltar: algo cambia de color o tamaño cuando toca mirarlo. Mover: una flecha o figura se desplaza para mostrar un recorrido. Las tres ayudan a entender; el texto que entra rebotando, no.}
\milcpaso{3}{milcTurquesa}{\textbf{La copla antes que las décimas.} El menú se escribe primero y lleva la esencia: si alguien solo ve el menú, ya sabe de qué trata todo. Las secciones desarrollan cada botón y vuelven, como cada décima vuelve a su verso.}
\milcpaso{4}{milcTurquesa}{\textbf{Cómo se hace, paso a paso.} (1) Cinco diapositivas con título. (2) Menú con cuatro botones. (3) Enlace de cada botón a su sección. (4) Contenido corto en cada sección. (5) Probar desde el menú. (6) Después, los regresos.}

\begin{drawbox}{milcTurquesa}{La presentación interactiva, en una mirada}
\textbf{Menú:} la pantalla índice, lleva la esencia. \\[1mm] \textbf{Sección:} una rama con su contenido. \\[1mm] \textbf{Enlace:} el botón que lleva de una pantalla a otra. \\[1mm] \textbf{Retorno:} el botón que vuelve al menú. \\[1mm] \textbf{Animación útil:} revelar, resaltar o mover para entender.
\end{drawbox}

\begin{softbox}{milcMostaza}{milcGris}{Errores comunes y su corrección}
(1) \textbf{Rama sin retorno}: quien entra a una sección no puede volver. (2) \textbf{Enlace roto}: el botón «glosario» no lleva a ninguna parte. (3) \textbf{Animación decorativa}: texto que rebota sin razón. (4) \textbf{Una transición distinta por pantalla}: cansa y distrae. (5) \textbf{Menú que no dice nada}: cuatro botones con «sección 1, 2, 3, 4» en vez del tema de cada una.
\end{softbox}

\begin{infoband}{milcTurquesa}


\textbf{Lo que va al cuaderno (Actividad 2 · CREA).} \textbf{Título:} «Actividad 2 --- El menú y las ramas». \textbf{Formato:} el texto de los cuatro botones del menú, la lista de las cuatro secciones con su contenido en una línea, y el enlace de la presentación. \textbf{Extensión:} media página. \textbf{Sabes que terminaste cuando} desde el menú entras a las cuatro secciones sin que ningún botón falle.
\end{infoband}

\puente{Ya tienes menú y ramas. Toca poner los regresos, las tres animaciones y probar que ningún camino deja atrapado a nadie.}

\milcestacion{milcMagenta}{Estación 3 de 3 · Hacer}

\begin{softbox}{milcMagenta}{milcGris}{Cómo vas a construir tu producto}
\textbf{Actividad 3 · EVALÚA --- Regresos, animaciones y la prueba del minuto} (25 min · individual). (1) Pon en cada sección un botón «volver al menú» enlazado a la primera diapositiva. (2) En la sección con estrella, subraya la palabra difícil y enlázala a una diapositiva nueva con su definición. Desde ahí pon un enlace de regreso a esa sección. (3) Agrega tres animaciones: una que revela, una que resalta y una que mueve. Para cada una, pregúntate: si la quito, ¿se entiende igual? Si sí, cámbiala. (4) Pásale el computador a tu pareja sin decir nada y cronometra un minuto: debe entrar a dos secciones y volver al menú. (5) Escribe la nota de cinco líneas: por qué estos caminos y no una secuencia.
\end{softbox}

\milcpaso{1}{milcMagenta}{\textbf{Tu entrega tiene tres piezas.} (1) La presentación compartida con enlace, con cinco pantallas, menú, retornos y glosario. (2) El resultado de la prueba del minuto. (3) La nota de cinco líneas sobre los caminos.}
\milcpaso{2}{milcMagenta}{\textbf{El retorno va en todas.} Cada sección y el glosario tienen botón de regreso. Es la regla de la décima: ninguna rama se queda suelta. Si tu pareja quedó atrapada en una pantalla, falta un retorno.}
\milcpaso{3}{milcMagenta}{\textbf{La prueba de quitar.} Una animación es útil solo si al quitarla se pierde algo. Revelar la respuesta después de la pregunta se pierde; un título que entra girando, no. Quita lo que no se pierde.}
\milcpaso{4}{milcMagenta}{\textbf{La prueba del minuto.} No expliques. Si tu pareja entró a dos secciones y volvió sola, los caminos funcionan. Si preguntó «¿y ahora cómo vuelvo?», ya sabes qué arreglar.}

\begin{drawbox}{milcMagenta}{Antes de entregar}
\checkbox Cinco pantallas, una es el menú con cuatro botones enlazados. \\[1mm] \checkbox Botón «volver al menú» en cada sección. \\[1mm] \checkbox Un glosario que lleva a la definición y devuelve. \\[1mm] \checkbox Tres animaciones que pasan la prueba de quitar. \\[1mm] \checkbox Mi pareja entró a dos secciones y volvió al menú en un minuto sin ayuda. \\[1mm] \checkbox Nota de cinco líneas sobre por qué estos caminos. \\[1mm] \checkbox Presentación compartida con enlace.
\end{drawbox}

\begin{softbox}{milcMagenta}{milcGris}{Plantilla de trabajo}
\textbf{Tema.} \textit{…} \\[1mm] \textbf{Menú.} \textit{Cuatro botones: …, …, …, …} \\[1mm] \textbf{Glosario.} \textit{La palabra … en la sección …} \\[1mm] \textbf{Animaciones.} \textit{Revela …; resalta …; mueve …} \\[1mm] \textbf{Prueba del minuto.} \textit{Mi pareja entró a … y … y volvió / se quedó en …}
\end{softbox}

\begin{infoband}{milcMagenta}


\textbf{Lo que va al cuaderno (Actividad 3 · EVALÚA).} \textbf{Título:} «Actividad 3 --- Regresos, animaciones y la prueba del minuto». \textbf{Formato:} la lista de las tres animaciones con qué se pierde si se quita cada una, el resultado de la prueba del minuto y la nota de cinco líneas. \textbf{Extensión:} una página. \textbf{Sabes que terminaste cuando} tu pareja volvió al menú sola y las tres animaciones tienen escrito qué se perdería sin ellas.
\end{infoband}

\puente{Lo que entregas hoy: la presentación con enlace y la nota de cinco líneas. La prueba de calidad: le pasas el computador a alguien, no le dices nada, y en un minuto entró a dos secciones y volvió al menú.}

\milcestacion{milcMostaza}{Producto de la sesión}
\textbf{Presentación de cinco pantallas con menú, retornos y glosario + tres animaciones útiles + prueba del minuto + nota de cinco líneas}.
\begin{softbox}{milcMostaza}{milcGris}{Criterios de calidad}
(1) El menú tiene cuatro botones que dicen el tema de cada sección y todos enlazan. (2) Cada sección y el glosario tienen botón de regreso que funciona. (3) Las tres animaciones pasan la prueba de quitar y hay escrito qué se pierde sin cada una. (4) Tu pareja entró a dos secciones y volvió al menú en un minuto sin ayuda. (5) La nota de cinco líneas explica por qué caminos y no secuencia. (6) La presentación está compartida con enlace.
\end{softbox}

\puente{Antes de cerrar, mira lo que hiciste desde cinco lados. Una presentación con cinco pantallas y regresos que funcionan vale más que una de veinte con animaciones que rebotan.}

\milcestacion{milcVino}{Evaluación liberadora · cinco dimensiones}

\begin{softbox}{milcVino}{milcGris}{Sentido de la evaluación}
Evaluar no es solo revisar una nota. Es reconocer qué aprendí, qué cambió en mí, qué emociones manejé, cómo me relaciono con la ciudad y la comunidad, y qué le dejaré a quien viene detrás.
\end{softbox}

\textbf{Mi reflexión liberadora en cinco dimensiones.} Completa cada frase iniciada:

\small
\begin{tabularx}{\textwidth}{>{\bfseries\RaggedRight\arraybackslash}p{3.4cm}Y>{\RaggedRight\arraybackslash}p{4.6cm}}
\rowcolor{milcVino}\color{white}Dimensión & \color{white}\bfseries Mi respuesta (frame starter) & \color{white}\bfseries Algo que descubrí\\
1. Personal & Lo que más cambió en mí fue \linead{6.5cm} & \linead{4.2cm}\\
2. Emocional & La emoción más fuerte fue \linead{2.5cm} y la regulé \linead{3cm} & \linead{4.2cm}\\
3. Ciudadana & En democracia esto importa porque \linead{6cm} & \linead{4.2cm}\\
4. Local (Cartago) & En mi barrio sería útil para \linead{6.5cm} & \linead{4.2cm}\\
5. Intergeneracional & A mi abuela/abuelo le diría \linead{6.5cm} & \linead{4.2cm}\\
\end{tabularx}
\normalsize

\begin{infoband}{milcVino}
\textbf{Para la próxima sustentación.} Vuelve a esta tabla en seis meses. Verás que las respuestas cambian — no porque lo aprendido cambie, sino porque tú cambias. La evaluación liberadora es una conversación contigo a través del tiempo.
\end{infoband}


\puente{Para terminar, tres ideas sobre dejar elegir el camino, contadas con palabras de esta guía: Dussel sobre el maestro que sabe oír, Séneca sobre el tiempo que derrochamos, y Floridi sobre saber preguntar.}

\milcestacion{milcGradeDark}{Tres ideas para cerrar · Dussel, un estoico y Floridi}

\begin{softbox}{milcVino}{milcGris}{Enrique Dussel · lente del nosotros}
\textit{El maestro que libera sabe oír en silencio a la juventud y al pueblo, en vez de imponer su orden.}

{\footnotesize\itshape\color{milcNegro!70}--- idea tomada de Enrique Dussel, contada con palabras de esta guía. No es una cita textual.}

Una presentación en línea recta impone el orden de quien la hizo. Un menú con caminos deja que quien mira elija por dónde entrar. Diseñar los regresos es una forma de oír al lector.

\textbf{¿Qué de lo que hice esta semana obligaba a los demás a ir en mi orden?}
\end{softbox}
\linead{\linewidth}\\[1mm]
\linead{\linewidth}

\begin{softbox}{milcMostaza}{milcGris}{Estoicismo · Séneca · cuidado interior}
\textit{La vida no es corta; nosotros la hacemos corta derrochando el tiempo.}

{\footnotesize\itshape\color{milcNegro!70}--- idea tomada de Séneca, contada con palabras de esta guía. No es una cita textual.}

Un título que entra rebotando le quita dos segundos a cada persona que lo mira, y no le da nada. Multiplícalo por veinte pantallas y treinta compañeros. La prueba de quitar cuida ese tiempo.

\textbf{¿Qué animación puse porque se veía bonita y no porque ayudara?}
\end{softbox}
\linead{\linewidth}\\[1mm]
\linead{\linewidth}

\begin{softbox}{milcTurquesa}{milcGris}{Luciano Floridi · lente de la infoesfera}
\textit{Ganan quienes saben preguntar y responder, y por eso saben qué merece recogerse y cuidarse.}

{\footnotesize\itshape\color{milcNegro!70}--- idea tomada de Luciano Floridi, contada con palabras de esta guía. No es una cita textual.}

Un buen menú es una lista de preguntas: ¿qué es?, ¿cómo funciona?, ¿para qué sirve? Cada sección responde una y vuelve. Un menú con «sección 1, 2, 3» no pregunta nada y nadie sabe por dónde entrar.

\textbf{¿Los botones de mi menú son preguntas que alguien se haría, o solo nombres?}
\end{softbox}
\linead{\linewidth}\\[1mm]
\linead{\linewidth}

\begin{infoband}{milcNegro}
\textbf{Nota para el docente.} Las tres ideas están contadas con palabras de esta guía a partir del pensamiento de cada autor; no son citas textuales. De 9.º en adelante se trabajan con la obra y su referencia.
\end{infoband}

\milcestacion{milcVino}{Mi autoevaluación · marca una casilla por fila}

{\small
\setlength{\extrarowheight}{2.2pt}
\begin{tabularx}{\textwidth}{>{\RaggedRight\arraybackslash\bfseries}p{3.0cm}YYY}
\rowcolor{milcVino}\color{white}\bfseries Criterio & \color{white}\bfseries Lo logré & \color{white}\bfseries En proceso & \color{white}\bfseries Necesito apoyo\\[.6mm]
Comprensión del tema & \milcopcion{Explico las ideas con mis palabras.} & \milcopcion{Reconozco las ideas, falta precisión.} & \milcopcion{Necesito repasar lo básico.}\\[.8mm]
\rowcolor{milcGris}Pensamiento computacional & \milcopcion{Usé los pilares al armar mi producto.} & \milcopcion{Usé algunos pilares.} & \milcopcion{Necesito releer la estación 2.}\\[.8mm]
Aplicación y producto & \milcopcion{Mi producto cumple los criterios.} & \milcopcion{Está armado, con ajustes pendientes.} & \milcopcion{Aún está en construcción.}\\[.8mm]
\rowcolor{milcGris}Reflexión 5 dimensiones & \milcopcion{Respondí cada una con honestidad.} & \milcopcion{Respondí algunas con detalle.} & \milcopcion{Mis respuestas son muy generales.}\\[.8mm]
Triángulo de pensamiento & \milcopcion{Conecté las 3 voces con mi trabajo.} & \milcopcion{Conecté al menos una.} & \milcopcion{No las relaciono con mi proceso.}\\[.8mm]
\end{tabularx}}

\vspace{3mm}
\textbf{Hoy entendí que:}\\[1mm]
\linead{\linewidth}\\[1mm]
\linead{\linewidth}

\textbf{Mi compromiso para la próxima sesión es:}\\[1mm]
\linead{\linewidth}\\[1mm]
\linead{\linewidth}

\begin{softbox}{milcMostaza}{milcGris}{Cita de despedida · Marco Aurelio}
\textit{``El obstáculo en el camino se vuelve el camino. Lo que estorba al camino se vuelve camino.''}\\[1mm]
Cada error de hoy, bien observado, se convierte en pieza del camino que pisarás mañana.
\end{softbox}

\small
\begin{tabularx}{\textwidth}{>{\bfseries}p{3.6cm}Y}
\rowcolor{milcGradeDark!12}Nombre completo: & \linead{10cm}\\
Fecha: & \linead{4cm} \quad Curso/grupo: \linead{4cm}\\
Mi firma: & \linead{9cm}\\
\end{tabularx}
\normalsize

\begin{infoband}{milcGradeDark}
\textbf{Para la próxima sesión.} Llega con tu presentación y la nota. En la sesión 3 vas a contar algo que te importa en tres minutos, con tu voz y sin leer. El menú que lleva la esencia es lo mismo que el gancho de un discurso.
\end{infoband}

\end{document}
